\section*{Abstract}
\addcontentsline{toc}{section}{Abstract}
\label{sec:abstract}

This project forms the environment sector of MODELBD, a cohort-wide data
integration initiative spanning seven sectors. Environmental information about
Bangladesh is published by several agencies in workbooks, reference files and
reports designed for human reading rather than integrated retrieval, so the
material stays scattered across separate organisations instead of one queryable
location. Differences in layout, naming, units, time grain, coverage and
missing-value notation make the publications difficult to compare directly. The
project addresses this problem by reviewing \mSourceDataFiles{} collected source
files from official and other dependable publishers, recording an exclusion
reason for material that carries no shared key or falls outside the approved
model, then selecting \mSourceFiles{} structured files and extracting
\mSourceBlocks{} traceable data blocks from four Bangladeshi organisations.
Reference registers and narrative publications are converted into structured
form where they define headings, units or station identity. Subject diagrams
drawn for individual datasets are consolidated into one integrated
entity--relationship model. The retained material is then reshaped, reconciled
and normalized to Boyce--Codd Normal Form without estimating unavailable source
values. The resulting SQLite database contains \mDbTables{} tables,
\mDbViews{} views, \mDbColumns{} columns and \mDbRows{} rows covering
\mYearMin--\mYearMax. Its declared relationships connect climate, time,
geography, rivers, water quality, forest and industrial records while preserving
their different daily, monthly, yearly and fiscal-period grains. Competency
queries, the final schema definition, primary- and foreign-key checks, a clean
rebuild and a backup--restore comparison evaluate the database; SQLite returns
\texttt{\mDbIntegrity} for the integrity check and \mDbFkViolations{}
foreign-key violations. The database therefore replaces manual reconciliation
across several spreadsheets with one query that returns the same answer to every
reader, which supports comparison, teaching and further environmental research.
The principal costs are the volume of manual cleaning the published layouts
demand, source values that disagree for the same key, uneven subject coverage
between years and the absence of any automatic refresh: \mProblemTotal{}
recorded problem occurrences span \mProblemClasses{} classes and
\mCellsMissingPct\% of the cells read are missing or unusable. Limitations stay
visible rather than hidden, particularly the \mUnmappedStations{} stations for
which the source material provides no supported district mapping. Later
revisions can widen the subject range, add further publishers and introduce a
scheduled reload.
